\section{Analisi Probabilistica e Algoritmi Randomizzati}

\begin{softnote}
	Finora abbiamo analizzato gli algoritmi nel caso peggiore.
	Questa sezione introduce un'alternativa: l'analisi del
	\textbf{caso medio}, e gli \textbf{algoritmi randomizzati}
	che rendono l'analisi media valida per qualunque input.
\end{softnote}

% ----------------------------------------------------------------
\subsection{Preliminari: variabili casuali indicatrici}

\begin{nota}
	Si assume che lo studente conosca il concetto di variabile casuale.
	Chi non lo conosce può studiare l'Appendice C di [CLRS] come
	riferimento sufficiente per questa sezione.
\end{nota}

\begin{definizione}[Variabile casuale indicatrice]
	Dato uno spazio campionario $S$ e un evento $A$, la
	\textbf{variabile casuale indicatrice} $I\{A\}$ è:
	\[
	I\{A\} = \begin{cases} 1 & \text{se } A \text{ si verifica} \\
		0 & \text{altrimenti} \end{cases}
	\]
\end{definizione}

\begin{lemma}
	Dato uno spazio campionario $S$ e un evento $A$, posto
	$X_A = I\{A\}$, si ha:
	\[
	E[X_A] = \Pr\{A\}.
	\]
\end{lemma}
\begin{proof}
	$E[X_A] = 1\cdot\Pr\{A\} + 0\cdot\Pr\{\bar{A}\} = \Pr\{A\}$.
\end{proof}

\begin{nota}
	\textbf{Perché sono utili?} Le v.c.\ indicatrici trasformano il
	calcolo di medie complesse in somme di probabilità semplici,
	grazie alla \textbf{linearità del valore atteso}:
	$E\!\left[\sum_i X_i\right] = \sum_i E[X_i]$,
	che vale anche quando le $X_i$ sono \textit{dipendenti}.
\end{nota}

\begin{esempio}
	\textbf{Numero atteso di teste in $n$ lanci di moneta.}
	Sia $X_i = I\{\text{testa all'$i$-esimo lancio}\}$,
	e $X = \sum_{i=1}^n X_i$.
	\[
	E[X] = \sum_{i=1}^n E[X_i] = \sum_{i=1}^n \frac{1}{2} = \frac{n}{2}.
	\]
	Notare che si evita il calcolo diretto della distribuzione binomiale.
\end{esempio}

% ----------------------------------------------------------------
\subsection{Analisi probabilistica: il Problema delle Assunzioni}

\begin{definizione}[Problema delle Assunzioni]
	\textbf{Input:} $n$ candidati ordinabili per bravura;
	costo di colloquio $c_c > 0$, costo di assunzione $c_a \gg c_c$.\\
	\textbf{Output:} assumere il candidato migliore con il minimo costo.\\
	\textbf{Strategia:} si assume il primo candidato; si sostituisce
	il corrente ogni volta che si incontra uno migliore.
\end{definizione}

Se $m$ è il numero di assunzioni effettuate, il costo totale è:
\[
n \cdot c_c + m \cdot c_a.
\]
Il costo $n\cdot c_c$ è fisso; quello che ci interessa è $m\cdot c_a$.
Nel caso peggiore $m = n$, nel caso migliore $m = 1$.
Vogliamo calcolare il \textbf{costo medio}, che dipende dalla
distribuzione dei candidati.

\paragraph{Calcolo del costo medio.}
Si assuma che le $n!$ permutazioni dei candidati siano equiprobabili.
Sia $X_i = I\{\text{il candidato $i$-esimo viene assunto}\}$,
$X = \sum_{i=1}^n X_i$.

Il candidato in posizione $i$ viene assunto se e solo se è il
migliore tra i primi $i$. Poiché ogni permutazione dei primi $i$
candidati è equiprobabile:
\[
\Pr\{X_i = 1\} = \frac{1}{i}.
\]

Quindi:
\[
E[X] = \sum_{i=1}^{n} E[X_i] = \sum_{i=1}^{n} \frac{1}{i}
\le \ln n + 1.
\]

\begin{teorema}[Costo medio del Problema delle Assunzioni]
	\[
	\text{Costo medio} \le c_a(\ln n + 1) = O(c_a \log n),
	\]
	contro il costo nel caso peggiore $O(c_a\, n)$.
\end{teorema}

\begin{nota}
	$\sum_{i=1}^n 1/i$ è la \textbf{serie armonica}. Vale l'approssimazione:
	$\ln(n+1) \le \sum_{i=1}^n \frac{1}{i} \le \ln n + 1.$
	Ad esempio, per $n=100$: $4.615 \le 5.177 \le 5.605$.
\end{nota}

% ----------------------------------------------------------------
\subsection{Algoritmi randomizzati}

\begin{softnote}
	Assumere una distribuzione uniforme dell'input non è sempre realistico.
	Gli algoritmi randomizzati risolvono questo problema imponendo
	\textbf{essi stessi} la distribuzione, indipendentemente dall'input.
\end{softnote}

\begin{center}
	\begin{tabular}{p{5.5cm}p{5.5cm}}
		\toprule
		\textbf{Analisi probabilistica} & \textbf{Algoritmo randomizzato} \\
		\midrule
		Algoritmo deterministico & Algoritmo non deterministico \\
		Stesso input $\to$ stesso risultato &
		Stesso input $\to$ esecuzione diversa \\
		Ipotesi sulla distribuzione dell'input &
		Nessuna ipotesi necessaria \\
		Caso peggiore dipende dall'input &
		Il comportamento non dipende da un input specifico \\
		\bottomrule
	\end{tabular}
\end{center}

\begin{osservazione}
	Nel Problema delle Assunzioni, l'algoritmo randomizzato permuta
	casualmente i candidati \textit{prima} di iniziare.
	In questo modo il costo medio $O(c_a \log n)$ è garantito
	\textbf{per qualunque input}, non solo per quelli ben distribuiti.
\end{osservazione}

% ----------------------------------------------------------------
\subsection{Permutazione casuale: \texttt{RandomizeInPlace}}
\begin{algorithm}[H]
	\caption{RandomizeInPlace$(A)$}
	\KwIn{Array $A[1:n]$}
	\KwOut{$A$ con una permutazione casuale uniforme}
	$n \gets |A|$\;
	\For{$i = 1$ \KwTo $n$}{
		scambia $A[i]$ con $A[\textsc{Random}(i, n)]$\;
	}
	\Return $A$\;
\end{algorithm}

\textbf{Complessità in tempo:} $O(n)$ — $n$ iterazioni, ciascuna con un
numero costante di operazioni e una chiamata a $\textsc{Random}(i,n)$,
assunta a costo costante.

\textbf{Complessità in spazio:} $O(1)$ — nessun array ausiliario.

\begin{lemma}
	\texttt{RandomizeInPlace} genera una permutazione casuale uniforme
	(ogni delle $n!$ permutazioni ha probabilità $1/n!$).
\end{lemma}
\begin{proof}
	Si dimostra tramite invariante di ciclo.
	
	\textbf{Invariante:} prima della $i$-esima iterazione, per ogni
	possibile $(i-1)$-permutazione, il sottoarray $A[1:i-1]$ la
	contiene con probabilità $\frac{(n-i+1)!}{n!}$.
	
	\textbf{Base} ($i=1$): il sottoarray $A[1:0]$ è vuoto;
	l'unica 0-permutazione (vuota) vi compare con probabilità
	$1 = n!/n!$. \checkmark
	
	\textbf{Conservazione:} si assume che ogni $(i-1)$-permutazione
	$\langle x_1,\ldots,x_{i-1}\rangle$ compaia in $A[1:i-1]$
	con probabilità $(n-i+1)!/n!$.
	La probabilità che l'algoritmo scelga $x_i$ come $i$-esimo
	elemento è $1/(n-i+1)$. Quindi la probabilità di ottenere
	$\langle x_1,\ldots,x_{i-1},x_i\rangle$ è:
	\[
	\frac{(n-i+1)!}{n!}\cdot\frac{1}{n-i+1}
	= \frac{(n-i)!}{n!}.
	\]
	Questo è esattamente l'invariante per $i+1$. \checkmark
	
	\textbf{Conclusione:} alla fine ($i=n+1$) ogni $n$-permutazione
	compare in $A[1:n]$ con probabilità
	$\frac{(n-n)!}{n!} = \frac{1}{n!}$.
\end{proof}
\begin{hardnote}
	
	\textbf{Da ricordare all'esame:} la correttezza di
	\texttt{RandomizeInPlace} si dimostra con un invariante di ciclo,
	non semplicemente contando le permutazioni. Il punto chiave è che
	a ogni passo la scelta uniforme in $A[i:n]$ garantisce che tutte
	le possibili estensioni siano equiprobabili.
\end{hardnote}

\subsection{QuickSort}

\begin{softnote}
	QuickSort è uno degli algoritmi di ordinamento più usati in pratica.
	A differenza di MergeSort, ordina \textbf{sul posto} e, pur avendo
	caso peggiore $\Theta(n^2)$, nel caso medio è $\Theta(n \log n)$
	con costanti moltiplicative molto piccole.
\end{softnote}

% ----------------------------------------------------------------
\subsection{Idea Generale}

QuickSort è un algoritmo che:
\begin{itemize}
	\item ordina \textbf{sul posto} come InsertionSort;
	\item è basato sulla tecnica \textbf{divide et impera} come MergeSort;
	\item ha complessità $\Theta(n^2)$ nel caso peggiore e $\Theta(n \log n)$
	nel caso medio (qui si dimostra il caso medio con l'ulteriore
	assunzione che i valori siano tutti distinti);
	\item nel caso medio le costanti moltiplicative nascoste da $\Theta$ sono
	abbastanza piccole $\Rightarrow$ molto efficiente in pratica, tra
	i più usati per l'ordinamento interno.
\end{itemize}

\begin{nota}
	In Java viene usato, nella sua versione parallela, quando si devono
	ordinare $\ge 44$ elementi.
\end{nota}

\paragraph{Schema divide et impera.}
\begin{itemize}
	\item \textbf{Divide:} partiziona l'array $A[p:r]$ in due
	sotto-array $A[p:q-1]$ e $A[q+1:r]$ tali che ogni elemento
	del primo sia $\le A[q]$ (il \emph{pivot}) e ogni elemento
	del secondo sia $> A[q]$.
	\item \textbf{Impera:} ordina i due sotto-array in modo ricorsivo
	fino a quando la dimensione è 0 o 1 (quando $p \ge r$).
	\item \textbf{Combina:} dato che i due sotto-array sono già ordinati
	al ritorno della chiamata ricorsiva, non serve alcuna ulteriore
	fase di combinazione.
\end{itemize}

\begin{algorithm}[H]
	\caption{QuickSort$(A, p, r)$}
	\KwIn{Array $A$, indici $p$, $r$}
	\KwOut{$A[p:r]$ ordinato in loco}
	\If{$p < r$}{
		$q \gets \text{Partition}(A, p, r)$\;
		$\text{QuickSort}(A, p, q-1)$\;
		$\text{QuickSort}(A, q+1, r)$\;
	}
\end{algorithm}

\begin{nota}
	Per ordinare un intero array $A$ di $n$ elementi si invoca
	\texttt{QuickSort}$(A, 1, n)$.
\end{nota}

% ----------------------------------------------------------------
\subsection{Partition$(A, p, r)$}

\begin{softnote}
	\texttt{Partition} è la componente fondamentale di QuickSort.
	Sceglie un valore \emph{pivot} (l'ultimo elemento $A[r]$),
	riorganizza l'array e restituisce l'indice $q$ della posizione
	finale del pivot.
\end{softnote}

\begin{algorithm}[H]
	\caption{Partition$(A, p, r)$}
	\KwIn{Array $A$, indici $p$, $r$}
	\KwOut{$q$ e $A$ modificato con $A[p:q-1] \le A[q] < A[q+1:r]$}
	$x \gets A[r]$\tcp*{Pivot: scegliere l'ultimo valore è utile per il \texttt{for}}
	$i \gets p - 1$\tcp*{Indice ultimo elemento sotto-array sinistro}
	\For{$j \gets p$ \KwTo $r-1$}{
		\If{$A[j] \le x$\tcp*{Se l'elemento deve andare nel sotto-array sx}}{
			$i \gets i + 1$\tcp*{Posizione per l'elemento da spostare}
			scambia $A[i]$ con $A[j]$\;
		}
	}
	$q \gets i + 1$\tcp*{Posizione per il pivot}
	scambia $A[q]$ con $A[r]$\;
	\Return $q$\;
\end{algorithm}

\subsubsection{Invariante di ciclo di Partition}

Per dimostrare la correttezza di \texttt{Partition}$(A, p, r)$,
si considera la seguente invariante di ciclo:

\begin{definizione}[Invariante di Partition]
	All'inizio di ogni iterazione del ciclo \texttt{for}, per qualsiasi
	indice $k$ dell'array vale almeno una delle seguenti proprietà:
	\begin{enumerate}
		\item Se $p \le k \le i$, allora $A[k] \le x$;
		\item Se $i+1 \le k \le j-1$, allora $A[k] > x$;
		\item Se $k = r$, allora $A[k] = x$ (il pivot).
	\end{enumerate}
\end{definizione}

\begin{nota}
	Non si considerano i valori con indice $t.c.\ j \le k \le r-1$
	perché sono valori non ancora considerati e quindi non ancora
	posizionati correttamente.
\end{nota}

\begin{proof}
	\begin{enumerate}
		\item \textbf{Inizializzazione} (prima della I iterazione, $i=p-1$,
		$j=p$): non ci sono valori tra $p$ e $i+1 = p$ né tra $i+1=p$
		e $j-1=p-1$. Le prime due proprietà dell'invariante sono
		banalmente soddisfatte. L'istruzione $x \gets A[r]$ soddisfa
		la III proprietà. \checkmark
		
		\item \textbf{Conservazione:} ci sono due possibilità.
		\begin{itemize}
			\item Se $A[j] > x$ (riga 4 non eseguita), allora non si fa
			nulla e $j' = j+1$: la proprietà 2 dell'invariante è
			soddisfatta all'inizio del prossimo giro perché $j' - 1 = j$.
			\item Se invece $A[j] \le x$, allora $i' = i+1$, vengono
			scambiati $A[i']$ e $A[j]$, poi $j' = j+1$. In seguito
			allo scambio, $A[i'] \le x$ e la proprietà 1 è soddisfatta.
			Anche $A[j'+1] > x$ perché era l'elemento $A[i]$ prima
			dello scambio. Quindi anche la proprietà 2 è soddisfatta.
		\end{itemize}
		
		\item \textbf{Conclusione:} alla fine del ciclo $j = r$. Ogni
		posizione dell'array si trova in uno dei tre insiemi descritti
		nell'invariante. L'istruzione alla riga 8 posiziona l'elemento
		pivot nella giusta posizione: \texttt{Partition} fa esattamente
		quello che ci si aspetta, riordinando l'array rispetto
		all'elemento pivot.
	\end{enumerate}
\end{proof}

\paragraph{Complessità di Partition.}
Il tempo di computazione ha sempre lo stesso ordine di grandezza
per qualsiasi input di dimensione $n = r - p + 1$:
\[
T_{\text{Partition}}(n) = \Theta(n).
\]

% ----------------------------------------------------------------
\subsection{Analisi di Complessità di QuickSort}

\begin{softnote}
	Il tempo di esecuzione di QuickSort dipende, a parità di $n$,
	da quanto il partizionamento è bilanciato.
\end{softnote}

\subsubsection{Caso peggiore}

\paragraph{Analisi elementare.}
\texttt{Partition} restituisce sempre un sotto-array di dimensione
$n-1$ e uno di dimensione $0$. In questo caso:
\[
T(n) = T(n-1) + T(0) + \Theta(n) = T(n-1) + \Theta(n),
\]
che è la somma aritmetica dei costi di \texttt{Partition}:
\[
T(n) = \Theta(n^2). \quad \text{(caso peggiore)}
\]

\begin{hardnote}
	Questo caso si verifica anche quando l'array è \textbf{completamente
		ordinato}! QuickSort è peggiore di InsertionSort in questo scenario.
\end{hardnote}

\paragraph{Analisi formale.}
\texttt{Partition} restituisce sempre due sotto-array, uno di dimensione
$d$ e l'altro di dimensione $n-d-1$, dove $0 \le d \le n-1$:
\[
T(n) = \max_{0 \le d \le n-1}\bigl[T(d) + T(n-d-1)\bigr] + \Theta(n).
\]
Si usa il metodo di sostituzione ipotizzando $T(n) \le cn^2$:
\[
T(n) \le \max_{0 \le d \le n-1}\bigl[cd^2 + c(n-d-1)^2\bigr] + \Theta(n).
\]
La quantità $d^2 + (n-d-1)^2$ ha il suo massimo quando $d=0$ o
$d=n-1$:
\[
\max_{0 \le d \le n-1}\bigl[d^2 + (n-d-1)^2\bigr] = (n-1)^2
\le n^2 - 2n + 1.
\]
Quindi $T(n) \le c(n^2 - 2n + 1) + \Theta(n) \le cn^2$, confermando
$T(n) = O(n^2)$.

\subsubsection{Caso migliore}

\texttt{Partition} restituisce sempre due sotto-array con dimensione
$\lfloor n/2 \rfloor$ e $\lceil n/2 \rceil - 1$ rispettivamente.
In questo caso (omettendo troncamenti e arrotondamenti):
\[
T(n) = 2T(n/2) + \Theta(n) \;\Rightarrow\; T(n) = \Theta(n \log n).
\quad \text{(caso migliore)}
\]

\subsubsection{Verso il caso medio}

\paragraph{Caso con proporzionalità fissa.}
Si supponga che \texttt{Partition} restituisca sempre due sotto-array
con dimensione $\frac{9}{10}n$ e $\frac{1}{10}n$ rispettivamente
(caso sbilanciato):
\[
T(n) = T\!\left(\tfrac{9}{10}n\right) + T\!\left(\tfrac{1}{10}n\right)
+ \Theta(n).
\]

\begin{softnote}
	Sembrerebbe un caso simile al peggiore perché sbilanciato... ma
	analizzando l'albero di ricorsione si vede che ogni livello ha
	costo $\Theta(n)$ finché il ramo $T(\frac{1}{10}n)$ non
	raggiunge la dimensione 1 (dopo $\log_{10} n$ livelli). Dopo tale
	livello ogni livello costa ancora $\Theta(n)$. La ricorsione termina
	alla profondità $\log_{10/9} n = \Theta(\log n)$, quindi il costo
	totale è $\Theta(n \log n)$.
\end{softnote}

\paragraph{Analisi verso il caso medio generale.}
In generale non si può supporre che tutte le chiamate ricorsive
siano sempre caso migliore, peggiore o con proporzionalità fissa.
Assumendo che gli input siano distribuiti in modo equiprobabile,
durante una computazione alcune ricorsioni divideranno bene,
altre no.

Si definisce l'elemento pivot $A[q]$ trovato da \texttt{Partition}
\textbf{buono} se $A[q]$ risiede tra il 25° e il 75° percentile:
\[
A[q] \text{ buono} \;\Leftrightarrow\;
|A[p:q-1]|, |A[q+1:r]| \le \tfrac{3}{4}|A[p:r]|.
\]
Il 50\% degli elementi sono buoni. Sia $X_q = I\{A[q]
\text{ è un elemento buono}\}$:
\[
\Pr\{X_q = 1\} = E[X_q] = \frac{1}{2}.
\]
Dati due pivot $A[q]$ e $A[q']$ scelti casualmente uno dopo l'altro,
il valore atteso del numero di pivot buoni nella sequenza è:
\[
E[X_q + X_{q'}] = E[X_q] + E[X_{q'}] = \frac{1}{2} + \frac{1}{2} = 1.
\]

Quindi, mediamente ogni 2 scelte casuali del pivot si ottiene
un valore buono. Due livelli consecutivi della ricorsione costano
complessivamente $\Theta(n)$. Dopo una partizione buona, la
dimensione del sotto-problema più grande è al più $\frac{3}{4}n$.
Raggruppando i livelli due a due, la ricorrenza si comporta come:
\[
T(n) = T\!\left(\tfrac{3}{4}n\right) + T\!\left(\tfrac{1}{4}n\right)
+ \Theta(n) \;\Rightarrow\; T(n) = \Theta(n \log n).
\]

% ----------------------------------------------------------------
\subsection{QuickSort Randomizzato}

\begin{softnote}
	Assumere che l'input abbia gli elementi distribuiti in modo casuale
	uniforme è troppo limitante. È preferibile mescolare l'input usando
	un generatore di numeri casuali, oppure scegliere il pivot
	casualmente prima di invocare \texttt{Partition}.
\end{softnote}

\begin{algorithm}[H]
	\caption{RandQuickSort$(A, p, r)$}
	\KwIn{Array $A$, indici $p$, $r$}
	\KwOut{$A[p:r]$ ordinato in loco}
	\If{$p < r$}{
		$q \gets \text{RandPartition}(A, p, r)$\;
		$\text{RandQuickSort}(A, p, q-1)$\;
		$\text{RandQuickSort}(A, q+1, r)$\;
	}
\end{algorithm}

\begin{algorithm}[H]
	\caption{RandPartition$(A, p, r)$}
	\KwIn{Array $A$, indici $p$, $r$}
	\KwOut{$q$ e $A$ modificato}
	$i \gets \textsc{Random}(p, r)$\;
	scambia $A[r]$ con $A[i]$\;
	\Return $\text{Partition}(A, p, r)$\;
\end{algorithm}

% ----------------------------------------------------------------
\subsection{Analisi del Caso Medio di QuickSort}

\begin{softnote}
	\texttt{RandQuickSort} garantisce una scelta uniforme del pivot ad ogni
	chiamata ricorsiva. Di seguito si presenta l'analisi formale
	della complessità nel caso medio seguendo [CLRS cap. 7].
\end{softnote}

Il costo reale dell'algoritmo è proporzionale al numero totale di
confronti (istruzione 4) eseguiti dalla procedura \texttt{Partition}.
Si deve quindi determinare il numero totale di confronti fatti da
\texttt{Partition} durante tutta una computazione di QuickSort.

Dato $A$, si rinominano tutti gli elementi come $z_1, z_2, \ldots, z_n$
dove $z_i$ è l'$i$-esimo elemento più piccolo di $A$.
QuickSort confronta in \texttt{Partition} una coppia $z_i, z_j$
\textbf{al più una volta} in tutta la sua esecuzione. Se uno dei
due diventa pivot, può essere confrontato con l'altro; un elemento
pivot non viene poi più confrontato con elementi appartenenti a
chiamate ricorsive successive.

Si definisce la variabile casuale indicatrice:
\[
X_{ij} = I\{z_i \text{ confrontato con } z_j\}.
\]

Il numero totale di confronti è:
\[
X = \sum_{i=1}^{n-1}\sum_{j=i+1}^{n} X_{ij}.
\]

Il numero totale atteso di confronti:
\[
E[X] = \sum_{i=1}^{n-1}\sum_{j=i+1}^{n} \Pr\{z_i
\text{ confrontato con } z_j\}.
\]

\paragraph{Calcolo di $\Pr\{z_i \text{ confrontato con } z_j\}$.}
Si consideri l'insieme $Z_{ij} = \{z_i, z_{i+1}, \ldots, z_j\}$.
\begin{itemize}
	\item Se \texttt{Partition} sceglie un elemento $x$ tale che
	$z_i < x < z_j$, allora $z_i$ e $z_j$ andranno su due
	sotto-array distinti e \textbf{non} verranno mai confrontati.
	\item Se $x = z_i$ oppure $x = z_j$, allora $z_i$ e $z_j$ verranno
	confrontati \textbf{una volta}.
\end{itemize}
Gli elementi di $Z_{ij}$ sono equi-probabili. La probabilità di
scegliere un elemento particolare è quindi $\frac{1}{j-i+1}$
(quando gli elementi sono distinti):
\[
\Pr\{z_i \text{ confrontato con } z_j\}
= \Pr\{z_i \text{ è il primo pivot di } Z_{ij}\}
+ \Pr\{z_j \text{ è il primo pivot di } Z_{ij}\}
= \frac{2}{j-i+1}.
\]

\paragraph{Calcolo finale.}
\begin{align*}
	E[X]
	&= \sum_{i=1}^{n-1}\sum_{j=i+1}^{n} \frac{2}{j-i+1} \\
	&= \sum_{i=1}^{n-1}\sum_{k=1}^{n-i} \frac{2}{k+1}
	\quad\text{(cambio variabile } k = j-i\text{)}\\
	&\le \sum_{i=1}^{n-1}\sum_{k=1}^{n} \frac{2}{k}
	= \sum_{i=1}^{n-1} O(\log n)
	= O(n \log n).
\end{align*}

\begin{teorema}[Complessità media di QuickSort]
	\[
	E[T(n)] = O(n \log n). \quad \text{(caso medio)}
	\]
\end{teorema}

\begin{center}
	\begin{tabular}{lll}
		\toprule
		\textbf{Caso} & \textbf{Input} & \textbf{Complessità} \\
		\midrule
		Peggiore & Array ordinato (o inversamente ordinato) & $\Theta(n^2)$ \\
		Migliore & Pivot sempre mediano & $\Theta(n \log n)$ \\
		Medio    & Input equiprobabili / pivot casuale       & $\Theta(n \log n)$ \\
		\bottomrule
	\end{tabular}
\end{center}

\begin{hardnote}
	\textbf{Da ricordare all'esame:} il caso peggiore di QuickSort è
	$\Theta(n^2)$ e si verifica su array \emph{già ordinati} con il pivot
	fisso all'ultimo elemento. La versione randomizzata
	\texttt{RandQuickSort} rende tale caso molto improbabile e garantisce
	$O(n \log n)$ in media per qualsiasi input.
\end{hardnote}

% ================================================================
